\noindent\textit{Proof.}
Fix a supported cell \(x\), and write \(c_x=P(C\mid X=x)\).  Proposition
\(\mathrm{prop{:}capacity\text{-}identification}\), whose premises include
conditional instrument independence, instrument overlap, and no defiers,
gives
\[
q_0(x)=P(S_0=1,C\mid X=x),
\qquad
q_1(x)=P(S_1=1,C\mid X=x).
\tag{1}
\]
If \(c_x=0\), both probabilities in (1) vanish because their events are
subsets of \(C\).  Every \(\ell_i(x)\) and \(h_j(x)\) also vanishes by
the same proposition.  Consequently
\[
q_0(x)=q_1(x)=m(x)=\Delta q(x)=0,
\quad
\Gamma_x^{\ast}=\Gamma_x^{0}=\{0\},
\quad
B_L(x)=B_U(x)=0.
\tag{2}
\]
There is no normalized survivor-complier distribution in this zero-mass
cell, so division below is restricted to \(c_x>0\).

Suppose \(c_x>0\).  The no-selection restriction states
\[
P(S_0=S_1=1\mid C,X=x)=1.
\tag{3}
\]
Since the event in (3) is contained in \(\{S_d=1\}\), the probability-one
upper bound gives \(P(S_d=1\mid C,X=x)=1\) for \(d=0,1\).  Multiplication
by \(c_x>0\), followed by (1), yields
\[
q_0(x)=q_1(x)=c_x,
\qquad \Delta q(x)=0,
\qquad m(x)=c_x>0.
\tag{4}
\]
Proposition \(\mathrm{prop{:}tie\text{-}face\text{-}collapse}\) therefore
implies \(\Gamma_x=\Gamma_x^{0}\).  Directly, if
\(\gamma_x\in\Gamma_x=\Gamma_x^{\ast}\), its row and column sums satisfy
\[
0\le r_i:=\sum_j\gamma_{ij,x}\le\ell_i(x),
\qquad
0\le s_j:=\sum_i\gamma_{ij,x}\le h_j(x).
\]
Exact mass and (4) imply
\[
\sum_i\{\ell_i(x)-r_i\}=q_0(x)-m(x)=0,
\qquad
\sum_j\{h_j(x)-s_j\}=q_1(x)-m(x)=0.
\]
All summands are nonnegative, so both margins bind coordinatewise.  This
proves \(\Gamma_x\subseteq\Gamma_x^{0}\); the reverse inclusion follows
because an exact-marginal coupling has total mass \(m(x)\).  Thus
\[
\Gamma_x=\Gamma_x^{\ast}=\Gamma_x^{0}.
\tag{5}
\]

Define the normalized cumulative distribution functions
\[
F_0(t)=\frac{\ell_{\le t}(x)}{m(x)},
\qquad
F_1(t)=\frac{h_{\le t}(x)}{m(x)},
\qquad F_0(-1)=0.
\tag{6}
\]
The capacity totals in (4) show that \(\ell_i(x)/m(x)\) and
\(h_j(x)/m(x)\) are probability mass functions.  The scaling map
\(\gamma_x\mapsto\gamma_x/m(x)\) is a bijection from \(\Gamma_x^{0}\)
onto their complete coupling polytope.

Definition \(\mathrm{def{:}threshold\text{-}cuts}\) and \(\Delta q(x)=0\)
give
\[
\frac{B_L(x)}{m(x)}
=\max\left\{0,\max_{t\in\mathcal Y}
\frac{\ell_{\le t}(x)-h_{\le t}(x)}{m(x)}\right\}
=\max\left\{0,\max_t\{F_0(t)-F_1(t)\}\right\}.
\tag{7}
\]
For the upper endpoint,
\(\ell_{<t}(x)/m(x)=F_0(t-1)\) and
\(h_{>t}(x)/m(x)=1-F_1(t)\), whence
\[
\frac{B_U(x)}{m(x)}
=\min\left\{1,\min_{t\in\mathcal Y}
\{F_0(t-1)+1-F_1(t)\}\right\}.
\tag{8}
\]
At \(t=0\), the inner expression is \(h_{>0}(x)/m(x)\le1\), so the
separate cap is redundant.  The revalidated cited Lemma
\(\mathrm{lem{:}ordinal\text{-}fixed\text{-}marginal\text{-}strict\text{-}benefit}\),
which records Lu, Ding, and Dasgupta's Proposition 2, equation (6), states
that (7)--(8) are exactly the sharp lower and upper probabilities of
\(Y_1>Y_0\) over all couplings of the normalized ordinal marginals.  The
scaling bijection and (5) prove the claimed cellwise reduction. \(\square\)